\section{Discussion}\label{sec:discussion}

\subsection{What a telemetry detector is for}
The finite-key audit splits the attacks cleanly. Attacks inside the device model never produced a breach; the decoy-state bound shortened or aborted the key on its own, as Remark~\ref{rem:inmodel} predicts. Attacks outside the model breached every block that produced a key. A telemetry detector therefore adds secrecy only against model-violating attacks, and Proposition~\ref{prop:indist} limits even that role to attackers who leave a statistical trace. Detection accuracy on intercept-resend or PNS, the headline numbers of several ML studies \cite{isha2026,alkuwari2026}, measures something the protocol already provides. Its operational value is real but different: faster incident response, and a signal that separates a denial-of-service from a degrading fibre. Reporting it as a security gain overstates it.

\subsection{Why the learned detectors fail under drift}
Drift is not noise around a fixed point. Misalignment grows for an hour and resets; loss follows the day. A model trained on stationary data learns a tight region around the calibration point, and the healthy link leaves that region within minutes. Domain randomisation widens the region but blurs the boundary to the attack classes, so RF-DR still raised false alarms at \fprRFDRDOne\% at $d=1$ with a seed-to-seed interval wider than the mean. IF-drift shows the other side: with representative benign drift it stays at \fprIFdriftDOne\% at $d=1$, within a small factor of the target, and fails once the drift leaves the range it was fitted on. The DCS avoids the problem by design rather than by data. It refits the three drifting parameters in every window, and only relations that drift cannot change enter the statistic. The cost is visible too. Whatever the refit can absorb, the statistic cannot see: pure loss, and the missing dark counts of Blind-nodark. The post-hoc DCS+Y0 variant shows that adding prior knowledge about one parameter recovers an attack and pays for it in false alarms when that parameter drifts.

\subsection{The limits of the indistinguishability argument}
Proposition~\ref{prop:indist} is an existence statement about one attacker, and it assumes a photon-number-resolving detector for Eve, knowledge of the current channel and the ability to blind Bob's detectors without side effects. Each assumption is demanding in practice, and each has been met in some form: blinding of commercial detectors \cite{lydersen2010} and a full-field eavesdropper that left the QBER unchanged \cite{gerhardt2011}. The static variant shows that partial knowledge already creates a trace, though in an unexpected direction: an attacker who replays a stationary channel is caught by the absence of drift, not by its presence. A defender can exploit this only with a detector that alarms on low variation, and a capable attacker removes the trace by tracking the channel.

The practical conclusion is not new, but the benchmark makes it quantitative: against detector-control attacks, telemetry is a weak line of defence. Hardware measures such as photocurrent monitoring \cite{yuan2010}, randomised detector efficiency \cite{lim2015}, with the caveats found in testing \cite{huang2016}, or measurement-device-independent protocols \cite{lo2012} change the observation model itself, and only they can turn Blind-cal into a detectable attack. A telemetry monitor should be judged by its false-alarm rate on a drifting link and by the model-violating attacks it catches, with an explicit statement of the attacker it cannot catch.

\subsection{Threats to validity}
\emph{Internal.} All results come from one simulator. The device model (\ref{eq:honest}) omits afterpulsing, detector dead time, efficiency mismatch and basis-dependent misalignment, so some benign effects that would raise DCS residuals in hardware are absent; the empirical-null standardisation is meant to absorb such effects, but its behaviour on real data is untested. Each attack is a model of the attack, not an implementation. The breach test (\ref{eq:breach}) is one-sided and counts only bits Eve knows exactly.

\emph{External.} The drift amplitudes are assumptions guided by qualitative reports \cite{dixon2010,dynes2012}, not fitted to a link; the scale $d$ and the sensitivity runs bracket them but cannot replace field data. The protocol is three-intensity decoy-state BB84 with one finite-key analysis \cite{lim2014}; other protocols and tighter analyses \cite{tomamichel2012} would change key rates but not the indistinguishability argument, which depends only on the photon-number structure.

\emph{Construct.} Per-window FPR at a fixed threshold is one operating point. Sequential procedures such as CUSUM \cite{page1954} trade delay against false alarms differently and could raise power for small intercept-resend fractions; the comparison here is window-level for all detectors alike. The ML baselines follow the feature sets and training styles of the cited work but are not reimplementations of those papers.

\emph{Statistical.} Ten seeds bound the attainable Wilcoxon $p$-values, so effect sizes are reported as intervals. The AUC near 0.5 for Blind-cal is an estimate with sampling spread (\aucCalMin--\aucCalMax), consistent with the exact equality that Proposition~\ref{prop:indist} proves.
